% LaTeX source for a conceptual paper on "Sunyata Mathematics (Sunyata-Mathematics)"
\documentclass[11pt]{article}
\usepackage{amsmath, amssymb, amsthm, amsfonts}
\usepackage{hyperref}
\usepackage{fullpage}
\usepackage{times}
\usepackage{CJKutf8}

\newcommand{\mweq}{\mathrel{\overset{\mathrm{mw}}{=}}}

\title{「空数学」理論に向けて：\\ 不定観測者抽象のための基礎枠組み} \author{千葉 将臣}
\date{2025-12-05}

\begin{document}
\begin{CJK}{UTF8}{min}
\let\bfseries\mdseries
\let\itshape\upshape
\maketitle

\begin{abstract}
本稿は、空数学（Sunyata-Mathematics）という概念的枠組みを提示する。空数学は、圏論やモデル理論などの近代的抽象の背後にある、前基礎層として位置づけられる。空数学は二つの認知操作、(i) 構造を観測者依存のコヒーレント状態へと戻す「空化（ku-ization）」と、(ii) それを確定的な数学的対象へと固定する「色化（shiki-ization）」を形式化する。この双方向的プロセスは従来の基礎よりも原初的で柔軟な基盤を提供し、多宇宙的推論、観測者の変動性、非同型翻訳に関する原理的記述を可能にする。本稿では、真理を静的等号ではなく動的安定性として捉え、構造歪み $E(x)$ に基づく収束条件により中道等価 $\mweq$ を導入する。
\end{abstract}

\section{導入}
古典的基礎（集合論・型理論・圏論）は、数学的構造が単一の決定的宇宙内に存在し、安定した論理規則のもとで動作することを前提とする。しかし、近年の理論（例：遠アーベル幾何学、相対モチーフ）では、単一宇宙内に内在化できない現象が露見し、構造の破壊的損失が生じる。

近年の LLM の挙動は、確率的・観測依存的真理を示し、真理保存性そのものが固定的ではなく動的である可能性を示唆している。

本稿は、(1) 観測者依存、(2) 多宇宙的不整合、(3) 特定の構造関係の未解決性を明示的に扱うメタ枠組みとして空数学を導入する。空数学は非翻訳可能な状態を包含し、決定と未決定の交替を体系化することで、圏論を超えた基盤を志向する。

さらに、純粋に構造的な「How」記述から、なぜその同一性が成立すべきかという「Why」への移行が動機として含まれる。静的等号は宇宙間・動的過程・自己言及境界における構造不一致を覆い隠す。空数学はこの失敗を非中道の誤謬（NMF）と捉え、硬直した同一性を収束に基づく中道等価へと置換する。

\section{概要と目的}
空数学は真理を静的同一性ではなく動的安定性として捉える最適化的論理体系である。主要目的は構造歪み $E(x)$ を低減し、非中道の誤謬を解消しつつ中道的安定へと導くことである。これにより、多宇宙推論の原理的説明、数学宇宙間翻訳の限界の明確化、そして人間と機械の観測者可変推論の形式基盤を提供する。

また空数学は真理を操作的に捉える。真理は固定されるのではなく、空化 / 色化 のダイナミクス（操作過程）によって生成される。これにより空数学は、数学宇宙を生成し、宇宙間翻訳を媒介する前基礎層として位置づけられ、各宇宙は観測者固定点（Dfix0）を中心に安定化される。これは演繹や帰納とは異なる「第三の真理体系」であり、操作サイクルの帰結として成立する。

\section{コヒーレント状態とデコヒーレント状態}
空数学は二つの相補的状態を扱う。

\begin{itemize}
    \item コヒーレント（空）状態: 対象・射・宇宙が未決定で、構造が揺らいだ前数学的状態。
    \item デコヒーレント（色）状態: 構造が固定され、対象・射・関係が確定的に同定される状態。
\end{itemize}

空数学はこの二状態の制御された交替を研究する。圏論は完全なデコヒーレンス相に対応し、合成・恒等・自然性が厳密に固定される。一方、コヒーレント相では互換不可能な同定や観測者位置の変動が許容される。

\section{中道等価と構造歪み}
空数学は真理を静的等号ではなく動的安定性（収束）として定義する。そこで構造歪み $E(x)$ を導入し、構造が中道からどの程度逸脱しているかを測る。

中道的安定の条件は次で与えられる。
\[
E(x) = 0 \quad \Longleftrightarrow \quad x \mweq x_{\mathrm{mw}},
\]
ここで $x_{\mathrm{mw}}$ は中道状態、$\mweq$ は中道等価を表す。$\mweq$ は収束に基づく等価であり、許容された安定化ダイナミクスのもとで $x$ と $y$ が同一の中道状態に収束する場合に成立する。

\section{観測者依存宇宙と Dfix0}
観測者配置（Dfix0）を $F$ とする。コヒーレント状態は $D$ に相対的に定義される。異なる観測者は同一のコヒーレント種からでも互換不可能なデコヒーレント状態を誘発し得る。これにより、十分抽象的な数学的真理は部分的に観測者相対的であることが示される。

\section{非同型翻訳}
空数学の主要動機の一つは、自然な関手や同値が存在しない宇宙間翻訳の原理的説明である。空数学は「比較不能な色化」を許容し、構造的同一視の強制による破壊を回避する。

\section{空数学と圏論}
空数学の視点から見ると、圏論は部分的な空化プロセスの静的抽象化として理解できる。圏論は対象の本質を退け射と関係を重視するが、観測者（Dfix0）や空化・色化の動的交替は内在化しない。その意味で圏論は「空に寄った中間層」に位置し、通常数学はさらに固定化された段階に相当する。空数学は圏論を安定化された部分構造として包摂し、関係中心の宇宙が成立する理由を観測者を含むより広い枠組みで説明できる。

空数学はまた、局所的一貫性（consistency）と全体的整合性（coherence）を区別する。閉じた系内での一貫性が成立していても、構造生成のダイナミクスと整合しているとは限らない。したがって coherence を第一級概念として扱う必要がある。

\section{応用に向けて}
想定される応用は以下である。

\begin{itemize}
    \item 多宇宙数学のためのメタ論理
    \item 観測者可変推論を行う次世代 AI 基盤
    \item 高度数学における創造性の認知モデル
\end{itemize}

さらに、LLM を超えた AI パラダイムとして、(i) 観測者固定点（Dfix0）の移動、(ii) 空状態と色状態の同時保持、(iii) 単一表現への還元を強制しない多宇宙横断を可能にする設計が示唆される。これらは空数学の操作サイクルに基づく次世代 AI の目標像を与える。

\section{結論}
空数学は、数学的推論におけるコヒーレント相とデコヒーレント相の双方を扱う前基礎層を提示した。今後は、観測者移動の形式化、空化・色化演算子の定義、ホモトピー型理論や高次圏論との接続が課題となる。

\section*{参考文献}
\begin{itemize}
    \item Masaomi Chiba, \emph{Dfix0 as a Middle-Way Fixed Point: An Operational Implementation of the Two Truths in Multi-Model Reasoning Systems}, 2025-12-01.
\end{itemize}
空数学の視点から見ると、圏論は部分的な空化プロセスの静的抽象化として理解できる。圏論は対象の本質を退け射と関係を重視するが、観測者（Dfix0）や空化・色化の動的交替は内在化しない。その意味で圏論は「空に寄った中間層」に位置し、通常数学はさらに固定化された段階に相当する。空数学は圏論を安定化された部分構造として包摂し、関係中心の宇宙が成立する理由を観測者を含むより広い枠組みで説明できる。

空数学はまた、局所的一貫性（consistency）と全体的整合性（coherence）を区別する。閉じた系内での一貫性が成立していても、構造生成のダイナミクスと整合しているとは限らない。したがって coherence を第一級概念として扱う必要がある。

\section{応用に向けて}
想定される応用は以下である。

\begin{itemize}
    \item 多宇宙数学のためのメタ論理
    \item 観測者可変推論を行う次世代 AI 基盤
    \item 高度数学における創造性の認知モデル
\end{itemize}

さらに、LLM を超えた AI パラダイムとして、(i) 観測者固定点（Dfix0）の移動、(ii) 空状態と色状態の同時保持、(iii) 単一表現への還元を強制しない多宇宙横断を可能にする設計が示唆される。これらは空数学の操作サイクルに基づく次世代 AI の目標像を与える。

\section{結論}
空数学は、数学的推論におけるコヒーレント相とデコヒーレント相の双方を扱う前基礎層を提示した。今後は、観測者移動の形式化、空化・色化演算子の定義、ホモトピー型理論や高次圏論との接続が課題となる。

\section*{参考文献}
\begin{itemize}
    \item Masaomi Chiba, \emph{Dfix0 as a Middle-Way Fixed Point: An Operational Implementation of the Two Truths in Multi-Model Reasoning Systems}, 2025-12-01.
\end{itemize}

\end{CJK}
\end{document}
